\scanpage{44}
\mspage{131}

d'où il résulte que les hom. de chgt. de base
$(***)$ sont iso, et comme il s'agit de
faisceaux loc. const. et que $p$ est une
Hot-équivalence, il s'ensuit que $(**)$ est
iso, donc aussi $(*)$ sur les aboutissements,
qed.

\uline{Définition 1} Sous les conditions équiv. du th. 1,
on dit que $f$ est une Hot-fibration.

\uline{Proposition 1} Cette notion est stable par chgt. de
base (trivial), par composition (\unsure{formel|formelle} sur a)). Pour
que $f$ soit une Hot-fibration, il f. et s. que $f$
le soit.
\deleted{J'avais envie de mettre les conditions duales
a$''$, b$''$ de a$'$, b$'$, mais il n'y avait pas de place
pour \unsure{elles|aller} dans la démonstration.}

\inserted{\uline{Prop. 2} Soit $f:X\to S$ une Hot-fibration. Alors $s\mapsto\mathrm{hot}(X_s)$ est un système local sur $S$ à val. dans Hot. Donc les $X_s$ pour $s$}
% margin note written sideways
\inserted{\unsure{dans une|d'une} \unsure{$=$} composante connexe de $S$ sont Hot-isomorphes entre eux. (mais pas canoniquement, si cette comp. n'est pas 1-connexe).}

\uline{Théorème 2} \deleted{\unsure{Soit $f$}}
\diagram{\begin{tikzcd}[ampersand replacement=\&,row sep=small,column sep=small]
X^{\unsure{\flat|1}}\arrow[rr,"f"]\arrow[dr,"\unsure{p|p^\flat}"']\&\&X^{\unsure{\sharp|2}}\arrow[dl,"\unsure{p^\sharp|p'}"]\\
\&S\&
\end{tikzcd}}
un diagramme commutatif dans Cat, avec $p,p'$ des Hot-fibrations, et soit $S_0\subset S$ une partie de $S$ telle
que $S_0\to\pi_0(S)$ soit surjectif. Pour que $f$ soit
une \inserted{Hot-}équivalence, il f. et s. que les
$f_s:X_s\unsure{\text{---}|\to}Y_s$ $(s\in S_0)$ le soient.

a) La suffisance se prouve comme tantôt,
avec le critère de Artin-Mazur : \deleted{Il faut}

% archive sheet number: 44
